\chapter*{Conclusion générale}
\addcontentsline{toc}{chapter}{Conclusion générale}
\markboth{Conclusion générale}{}

Ce rapport a présenté les travaux réalisés durant le stage de perfectionnement effectué
au sein du département Informatique (IT) du Mövenpick Hôtel de Tanger, portant sur la
conception et le développement de \textbf{NetSentry}, un outil interne de découverte et
de cartographie des équipements connectés au réseau local de l'hôtel.

Le premier chapitre a permis de poser le contexte général du stage : l'organisme
d'accueil, son organisation interne et son environnement, ainsi que le département
Informatique et la mission confiée durant ces deux mois. Le deuxième chapitre a détaillé
la méthodologie de travail suivie -- analyse du besoin, cahier des charges et choix
technologiques -- qui a cadré l'ensemble du projet. Le troisième chapitre a présenté la
modélisation du système : architecture réseau, diagramme de cas d'utilisation, diagramme
de séquence, modèle conceptuel de données, ainsi que l'organisation du travail (planning
prévisionnel et suivi Kanban des tâches). Enfin, le quatrième chapitre a illustré, à
travers des captures d'écran, le fonctionnement réel de l'outil développé.

\section*{Apports du stage}

Les premières semaines ont permis d'acquérir des notions de réseau (adressage IP/MAC,
ports, protocole TCP/IP) ainsi qu'une prise en main de l'outil Nmap, nécessaires à la
réalisation du projet. Ce stage a également permis de mieux comprendre le fonctionnement
d'un département informatique en entreprise et la collaboration avec un prestataire
externe pour les systèmes métiers.

La principale difficulté rencontrée en début de stage a été de définir un sujet
réalisable compte tenu du périmètre limité des systèmes accessibles, les systèmes
métiers de l'hôtel étant gérés par un prestataire externe. Le sujet a ainsi été adapté
pour se concentrer sur le réseau interne, seul domaine réellement accessible au sein du
département IT.

\section*{Opinion sur l'entreprise et son fonctionnement}

Le Mövenpick Hôtel de Tanger est une structure bien organisée, faisant partie d'un
groupe hôtelier international reconnu. L'accueil au sein du département Informatique a
été positif, avec une réelle volonté de faire participer le stagiaire à des missions
concrètes. Un point d'amélioration identifié concerne la dépendance vis-à-vis d'un
prestataire externe pour la gestion des systèmes métiers, ce qui limite parfois la marge
de man\oe{}uvre du département IT interne sur certains sujets.

Ce qui a le plus plu au cours de ce stage est la possibilité de concevoir et développer
un outil répondant à un besoin réel du département, en combinant des compétences de
développement logiciel et des notions de sécurité réseau. La principale difficulté a
résidé dans le périmètre limité des systèmes accessibles, nécessitant d'adapter le
projet aux contraintes réelles de l'entreprise plutôt qu'à une idée initiale plus
ambitieuse.

\section*{Impact sur le projet professionnel}

Ce stage a permis de mieux cerner l'articulation entre développement logiciel et
infrastructure réseau, et de confirmer un intérêt pour les métiers à la croisée du
développement et de l'administration des systèmes et réseaux. Il constitue une première
expérience concrète du fonctionnement d'un département informatique en entreprise,
utile pour orienter la suite du parcours professionnel.
